\begin{thebibliography}{99}

\bibitem{abdelgawad2022}
Abdelgawad, A., et al. (2022).
\newblock An agent-based model of brain connectivity and alpha-synuclein spreading in Parkinson's disease.
\newblock \emph{Network Neuroscience}, 6(4), 1163--1187.

\bibitem{boelts2023}
J.~Boelts, J.-M.~Lueckmann, R.~Gao, and J.~H.~Macke, ``Simulation-based inference for efficient identification of generative models in computational connectomics,'' \textit{PLoS Comput.\ Biol.}, vol.~19, no.~9, e1011406, 2023, doi:10.1371/journal.pcbi.1011406.

\bibitem{brooks1990} D. J. Brooks, V. Ibanez, G. V. Sawle, et al., ``Differing patterns of striatal 18F-dopa uptake in Parkinson's disease, multiple system atrophy, and progressive supranuclear palsy,'' Annals of Neurology, vol. 28, no. 4, pp. 547--555, 1990. doi: 10.1002/ana.410280412

\bibitem{casolo2025}
C.~Casolo, et al., ``Identifiability challenges in sparse linear ordinary differential equations,'' \textit{arXiv preprint arXiv:2501.xxxxx}, 2025.

\bibitem{chu2019}
Y.~Chu, et al., ``Intrastriatal alpha-synuclein fibrils in monkeys: spreading, imaging and neuropathological changes,'' \textit{Brain}, vol.~142, no.~11, pp.~3565--3579, 2019. doi:~10.1093/brain/awz261.

\bibitem{conte2024}
M.~Conte, et al., ``Structural and practical identifiability of contrast transport models for DCE-MRI,'' \textit{PLoS Computational Biology}, vol.~20, no.~1, e1011773, 2024. doi:~10.1371/journal.pcbi.1011773.

\bibitem{dong2023}
Dong, R., et al. (2023).
\newblock Differential elimination for dynamical models via projections with applications to structural identifiability.
\newblock \emph{SIAM Journal on Applied Algebra and Geometry}, 7(1), 194--235.

\bibitem{dupre2026paper10} B.\,D.~Dupre, ``Bidirectional-Ready Mechanistic Twin for Parkinson's Disease: External Validation, Calibration-Aware Observation Likelihood, and NASEM 2024 Self-Audit,'' 2026 (submitted to \textit{npj Parkinson's Disease}).

\bibitem{dupre2026paper7} B.\,D.~Dupre, ``Per-Patient Bayesian Calibration of a Coupled $\alpha$-Synuclein Aggregation--Neuron Death ODE in Parkinson's Disease,'' 2026 (submitted to \textit{CPT: Pharmacometrics Syst.\ Pharmacol.}).

\bibitem{dupre2026paper8b} B.\,D.~Dupre, ``Regional Dopamine Transporter Decline Rates from Serial DaT-SPECT in Parkinson's Disease,'' 2026 (submitted to \textit{Mov.\ Disord.}).

\bibitem{dupre2026phase2}
Dupre, B. (2026).
\newblock Per-patient toxicity flux estimation from longitudinal DaT-SPECT in Parkinson's disease.
\newblock \emph{Manuscript in preparation}.

\bibitem{fearnley1991}
J.~M.~Fearnley and A.~J.~Lees, ``Ageing and Parkinson's disease: substantia nigra regional selectivity,'' \textit{Brain}, vol.~114, no.~5, pp.~2283--2301, 1991.

\bibitem{fiorenzato2021}
E.~Fiorenzato, et al., ``Asymmetric dopamine transporter loss affects cognitive and motor progression in Parkinson's disease,'' \textit{Movement Disorders}, vol.~36, no.~10, pp.~2303--2314, 2021. doi:~10.1002/mds.28790.

\bibitem{fornari2019}
S.~Fornari, A.~Sch{\"a}fer, M.~Jucker, and A.~Goriely, ``Prion-like spreading of Alzheimer's disease within the brain's connectome,'' \textit{J.\ R.\ Soc.\ Interface}, vol.~16, no.~159, Art.\ no.~20190356, 2019.

% --- G ---

\bibitem{garbarino2021}
S.~Garbarino, et al., ``Investigating hypotheses of neurodegeneration by learning dynamical systems of protein propagation in the brain,'' \textit{NeuroImage}, vol.~235, 118024, 2021. doi:~10.1016/j.neuroimage.2021.118024.

\bibitem{gutenkunst2007}
R.~N.~Gutenkunst, J.~J.~Waterfall, F.~P.~Casey, K.~S.~Brown, C.~R.~Myers, and J.~P.~Sethna, ``Universally sloppy parameter sensitivities in systems biology models,'' \textit{PLoS Computational Biology}, vol.~3, no.~10, p.~e189, 2007, doi:10.1371/journal.pcbi.0030189.

% --- H ---

\bibitem{hashemi2024}
M.~Hashemi, A.~N.~Vattikonda, V.~Jirsa, and M.~Woodman, ``Simulation-based inference on virtual brain models of disorders,'' \textit{Mach.\ Learn.: Sci.\ Technol.}, vol.~5, no.~3, 035019, 2024, doi:10.1088/2632-2153/ad6230.

\bibitem{kim2022}
H.K.~Kim, et al., ``Temporal trajectory model for dopaminergic input to the striatal subregions in Parkinson's disease,'' \textit{Parkinsonism \& Related Disorders}, vol.~105, pp.~40--46, 2022. doi:~10.1016/j.parkreldis.2022.08.022.

\bibitem{kish1988}
Kish, S.J., et al. (1988).
\newblock Uneven pattern of dopamine loss in the striatum of patients with idiopathic Parkinson's disease.
\newblock \emph{New England Journal of Medicine}, 318(14), 876--880.

\bibitem{lee2019}
C.~Lee, J.~Yoon, and M.~van~der~Schaar, ``Dynamic-DeepHit: A deep learning approach for dynamic survival analysis with competing risks based on longitudinal data,'' \textit{IEEE Trans.\ Biomed.\ Eng.}, vol.~67, no.~1, pp.~122--133, Jan. 2020.

\bibitem{modrak2022}
M.~Modrak, et al., ``Simulation-based calibration checking for Bayesian computation: the choice of test quantities shapes sensitivity,'' \textit{Bayesian Analysis}, vol.~1, no.~1, pp.~1--28, 2022. doi:~10.1214/22-BA1404.

\bibitem{pandya2019}
S.~Pandya, et al., ``Predictive model of spread of Parkinson's pathology using network diffusion,'' \textit{NeuroImage}, vol.~192, pp.~178--194, 2019. doi:~10.1016/j.neuroimage.2019.02.058.

\bibitem{pasquini2019}
J.~Pasquini, et al., ``Clinical implications of early caudate dysfunction in Parkinson's disease,'' \textit{Journal of Neurology, Neurosurgery, and Psychiatry}, vol.~90, no.~6, pp.~681--687, 2019. doi:~10.1136/jnnp-2018-320157.

\bibitem{patterson2019}
J.R.~Patterson, et al., ``Time course and magnitude of alpha-synuclein inclusion formation and nigrostriatal degeneration in the rat model of synucleinopathy triggered by intrastriatal alpha-synuclein preformed fibrils,'' \textit{Neurobiology of Disease}, vol.~121, pp.~1--12, 2019. doi:~10.1016/j.nbd.2018.09.001.

\bibitem{powell2018}
F.~Powell, M.~Tosun, R.~Sadeghi, P.~Weiner, and A.~Raj, ``Preserved structural network organization mediates pathology spread in Alzheimer's disease despite loss of white matter tract integrity,'' \textit{Journal of Alzheimer's Disease}, vol.~65, pp.~747--764, 2018. doi:~10.3233/JAD-180515.

% === Paper 8a identifiability additions (2026-04-11) ===

\bibitem{ppmi2011}
Parkinson Progression Marker Initiative (2011).
\newblock The Parkinson Progression Marker Initiative (PPMI).
\newblock \emph{Progress in Neurobiology}, 95(4), 629--635.

\bibitem{raj2012}
A.~Raj, A.~Kuceyeski, and M.~Weiner, ``A network diffusion model of disease progression in dementia,'' \textit{Neuron}, vol.~73, no.~6, pp.~1204--1215, 2012.

\bibitem{raue2009}
A.~Raue, C.~Kreutz, T.~Maiwald, J.~Bachmann, M.~Schilling, U.~Klingm{\"u}ller, and J.~Timmer, ``Structural and practical identifiability analysis of partially observed dynamical models by exploiting the profile likelihood,'' \textit{Bioinformatics}, vol.~25, no.~15, pp.~1923--1929, 2009, doi:10.1093/bioinformatics/btp358.

\bibitem{ren2025}
J.~Ren, L.~Bhatt, G.~Guo, and A.~Raj, ``Connectome-based biophysical models of pathological protein spreading in neurodegenerative diseases,'' \textit{PLoS Comput.\ Biol.}, vol.~21, no.~1, e1012743, 2025, doi:10.1371/journal.pcbi.1012743.

\bibitem{schafer2021}
A.~Schafer, et al., ``Bayesian physics-based modeling of tau propagation in Alzheimer's disease,'' \textit{Frontiers in Physiology}, vol.~12, 702975, 2021. doi:~10.3389/fphys.2021.702975.

\bibitem{simpson2024}
M.J.~Simpson, et al., ``Making predictions using poorly identified mathematical models,'' \textit{Bulletin of Mathematical Biology}, vol.~86, no.~7, 83, 2024. doi:~10.1007/s11538-024-01294-0.

\bibitem{sung2016}
C.~Sung, et al., ``Longitudinal decline of striatal subregional [18F]FP-CIT uptake in Parkinson's disease,'' \textit{Nuclear Medicine and Molecular Imaging}, vol.~50, no.~4, pp.~314--320, 2016. doi:~10.1007/s13139-016-0436-4.

\bibitem{talts2018}
S.~Talts, et al., ``Validating Bayesian inference algorithms with simulation-based calibration,'' \textit{arXiv preprint arXiv:1804.06788}, 2018.

% ============================================================================
% Phase 4 PK/PD defensive citations (added 2026-04-12)
% ============================================================================

\bibitem{tossicibolt2017}
L.~Tossici-Bolt \textit{et~al.}, ``[$^{123}$I]FP-CIT ENC-DAT normal database: the impact of the reconstruction and quantification methods,'' \textit{EJNMMI Phys.}, vol.~4, Art.\ no.~8, 2017.

\bibitem{villaverde2016}
A.~F.~Villaverde, A.~Barreiro, and A.~Papachristodoulou, ``Structural identifiability of dynamic systems biology models,'' \textit{PLoS Computational Biology}, vol.~12, no.~10, p.~e1005153, 2016, doi:10.1371/journal.pcbi.1005153.

\bibitem{vogel2023}
Vogel, J.W., et al. (2023).
\newblock Connectome-based modelling of neurodegenerative diseases: towards precision medicine and mechanistic insight.
\newblock \emph{Nature Reviews Neuroscience}, 24(10), 620--639.

\bibitem{wang2025}
S.~Wang \textit{et~al.}, ``Conditional neural ordinary differential equations for brain morphometry dynamics in Parkinson's disease,'' \textit{arXiv}, 2025; 2410.23429.

\bibitem{wieland2021}
F.-G.~Wieland, A.~L.~Hauber, M.~Rosenblatt, C.~T\"onsing, and J.~Timmer, ``On structural and practical identifiability,'' \textit{Current Opinion in Systems Biology}, vol.~25, pp.~60--69, 2021.

\bibitem{xu2024}
Xu, Y., et al. (2024).
\newblock Kinetic analysis of alpha-synuclein aggregation.
\newblock \emph{Nature Communications}, 15, 7083.

\bibitem{zheng2019}
Zheng, Y.Q., et al. (2019).
\newblock Local vulnerability and global connectivity jointly shape neurodegenerative disease propagation.
\newblock \emph{PLoS Biology}, 17(11), e3000357.

\end{thebibliography}
